\subsection*{الف}
نمادگذاری: $\mathcal{S}=\text{Sublayer}$، $\Gamma=\mathrm{diag}(\gamma)$، $\hat{x}=\dfrac{y-\mu\mathbf{1}}{\sigma}$ و $d$ بُعد ویژگی.

برای \lr{Post-LN}، با $y_l=x_l+\mathcal{S}(x_l)$ و $x_{l+1}=\mathrm{LN}(y_l)$:
\begin{flalign*}
\frac{\partial \mathcal{L}}{\partial x_l}
= \frac{\partial \mathcal{L}}{\partial x_{l+1}}\, J_{\mathrm{LN}}(y_l)\,\Big(I+\frac{\partial \mathcal{S}}{\partial x_l}\Big) &&
\end{flalign*}

ژاکوبین \lr{LayerNorm}:
\begin{flalign*}
J_{\mathrm{LN}}=\frac{\partial \mathrm{LN}(y)}{\partial y}
=\frac{\Gamma}{\sigma}\Big(I-\tfrac{1}{d}\mathbf{1}\mathbf{1}^\top-\tfrac{1}{d}\hat{x}\hat{x}^\top\Big) &&
\end{flalign*}

اثر اندازه: ضریب $\gamma/\sigma$ نُرم گرادیان را تغییر می‌دهد (با بزرگ‌شدن $\sigma$، گرادیان کوچک می‌شود). اثر جهت: عملگر تصویرِ $P=I-\tfrac1d\mathbf{1}\mathbf{1}^\top-\tfrac1d\hat{x}\hat{x}^\top$ مؤلفه‌های در راستای $\mathbf{1}$ (میانگین) و $\hat{x}$ (خودِ ورودیِ نرمال‌شده) را حذف می‌کند؛ پس $J_{\mathrm{LN}}\neq cI$، رتبه را کاهش و جهت گرادیان را تغییر می‌دهد.

در عمق $L$، گرادیانِ مسیر میان‌بر نیز در هر لایه در $J_{\mathrm{LN}}$ ضرب می‌شود:
\begin{flalign*}
\frac{\partial \mathcal{L}}{\partial x_0}
= \frac{\partial \mathcal{L}}{\partial x_L}\prod_{l=0}^{L-1} J_{\mathrm{LN}}^{(l)}\Big(I+\tfrac{\partial \mathcal{S}^{(l)}}{\partial x_l}\Big) &&
\end{flalign*}
حاصل‌ضربِ $L$ ژاکوبین با ضریب $\gamma/\sigma$ رشد یا افت هندسی (انفجار/محوشدگیِ گرادیان) ایجاد می‌کند.

برای \lr{Pre-LN}، با $x_{l+1}=x_l+\mathcal{S}(\mathrm{LN}(x_l))$:
\begin{flalign*}
\frac{\partial \mathcal{L}}{\partial x_l}
&= \frac{\partial \mathcal{L}}{\partial x_{l+1}}\Big(I+J_{\mathcal{S}}^{(l)}J_{\mathrm{LN}}^{(l)}\Big) \\
\frac{\partial \mathcal{L}}{\partial x_0}
&= \frac{\partial \mathcal{L}}{\partial x_L}\prod_{l=0}^{L-1}\Big(I+J_{\mathcal{S}}^{(l)}J_{\mathrm{LN}}^{(l)}\Big)
= \frac{\partial \mathcal{L}}{\partial x_L}\Big(I+\textstyle\sum_l J_{\mathcal{S}}^{(l)}J_{\mathrm{LN}}^{(l)}+\cdots\Big) &&
\end{flalign*}
جملهٔ پیشروِ $I$ یک مسیرِ همانیِ بدونِ نرمال‌سازی است؛ گرادیان بدون ضرب‌شدنِ مکرر در $J_{\mathrm{LN}}$ به همهٔ لایه‌ها می‌رسد، پس مقیاسِ گرادیان مستقل از عمق و کران‌دار می‌ماند.

نتیجه: نبودِ حاصل‌ضربِ هندسیِ ژاکوبین‌ها در \lr{Pre-LN}، گرادیان را در هر عمق پایدار نگه می‌دارد و آموزشِ مدل‌های عمیق پایدارتر می‌شود. در \lr{Post-LN} گرادیانِ لایه‌های نزدیک خروجی در آغاز بسیار بزرگ است و نرخ یادگیریِ زیاد، تخمین‌های وفقیِ \lr{Adam} را بی‌ثبات می‌کند؛ \lr{warmup} برای مهارِ همین گرادیانِ بزرگِ اولیه لازم است. چون \lr{Pre-LN} مقیاسِ گرادیان را از ابتدا کران‌دار و مستقل از عمق می‌سازد، نیاز به \lr{warmup} کاهش می‌یابد یا حذف می‌شود.

\subsection*{ب}
(۱) عمقِ مؤثر: در \lr{Pre-LN} به‌دلیلِ مسیرِ همانیِ غالب،
$x_L=x_0+\sum_l \mathcal{S}_l(\mathrm{LN}(x_l))$،
واریانسِ جریانِ باقی‌مانده با عمق انباشته می‌شود ($\mathrm{Var}(x_l)\sim l$). پس سهمِ نسبیِ هر زیرلایه
$\frac{\|\mathcal{S}_l\|}{\|x_l\|}\sim \frac{1}{\sqrt{l}}$
کوچک می‌شود و لایه‌های عمیق تقریباً همانی عمل می‌کنند؛ یعنی عمقِ مؤثرِ تبدیلِ ویژگی کم می‌شود و ظرفیتِ مدل کامل به‌کار نمی‌رود. در \lr{Post-LN} نرمال‌سازیِ پس از هر جمع، سهمِ همهٔ لایه‌ها را هم‌مقیاس نگه می‌دارد و تبدیلِ ویژگیِ غنی‌تری می‌دهد.

(۲) توزیعِ \lr{activation} ورودی به لایهٔ نهایی: آخرین عملیاتِ \lr{Post-LN} یک \lr{LayerNorm} است، پس ورودیِ سرِ خروجی نرمال (میانگین صفر، مقیاسِ کنترل‌شده) و خوش‌وضع است. در \lr{Pre-LN} آخرین عملیات یک جمع است و \lr{activation}های نهایی نرمال‌نشده و با واریانسِ بزرگ وارد سر می‌شوند؛ این شرطِ بدتر، بهینه‌سازیِ سر و برازشِ نهایی را اندکی ضعیف‌تر می‌کند.

(۳) نقشِ نرمال‌سازی به‌عنوانِ \lr{regularization}: در \lr{Post-LN} نرمال‌سازی روی مسیرِ اصلیِ سیگنال در هر لایه اعمال می‌شود و با محدودکردنِ مقیاس و توزیعِ بازنمایی‌های میانی یک \lr{regularization} ضمنی ایجاد می‌کند که در مقیاسِ کوچک تعمیم را بهبود می‌دهد. \lr{Pre-LN} فقط ورودیِ شاخه را نرمال می‌کند و جریانِ باقی‌مانده آزاد می‌ماند؛ کمبودِ این اثرِ تنظیمی می‌تواند به کمینهٔ کمی بدتر یا بیش‌برازش بینجامد.

جمع‌بندی: \lr{Pre-LN} پایداری را با هزینهٔ کاهشِ عمقِ مؤثر، بدترشدنِ وضعِ ورودیِ لایهٔ نهایی و کم‌شدنِ \lr{regularization} ضمنی به‌دست می‌آورد؛ در مدل‌های کوچک که پایداری گلوگاه نیست، این هزینه‌ها غالب می‌شوند و عملکردِ نهایی اندکی پایین‌تر می‌آید.

\subsection*{ج}
نمونه: \lr{Fixup initialization} و نسخهٔ ترنسفورمریِ آن \lr{T-Fixup}؛ روشِ جایگزینِ دیگر \lr{DeepNorm} است.

ایدهٔ \lr{Fixup}/\lr{T-Fixup}: حذفِ کاملِ \lr{LayerNorm} و کنترلِ مقیاس فقط از راهِ مقداردهیِ اولیه:
\begin{itemize}
\item آخرین لایهٔ وزنیِ هر شاخهٔ باقی‌مانده با صفر مقداردهی می‌شود تا در آغاز هر بلوک دقیقاً همانی باشد؛ پس واریانسِ سیگنالِ پیش‌رو در عمق حفظ می‌شود.
\item سایر وزن‌ها با ضریبِ $L^{-1/(2m-2)}$ ($m$ تعدادِ لایه‌های شاخه) کوچک می‌شوند تا تغییرِ خروجی پس از یک گامِ بهینه‌سازی به‌ازای هر بلوک $\Theta(1/L)$ و در کل $\Theta(1)$ بماند.
\end{itemize}
\lr{DeepNorm}: مسیرِ باقی‌مانده را بزرگ‌نمایی می‌کند، $x_{l+1}=\mathrm{LN}(\alpha x_l+\mathcal{S}(x_l))$ با $\alpha>1$، و برخی وزن‌ها را با $\beta<1$ کوچک می‌کند تا اندازهٔ به‌روزرسانیِ مدل مستقل از عمق کران‌دار شود.

توضیحِ نظری: کاری که نرمال‌سازی انجام می‌دهد، تأمینِ سه نامتغیّرِ مقیاسی در سراسرِ عمق است:
\begin{flalign*}
\text{(۱) }\mathrm{Var}(x_l)=\Theta(1),\qquad
\text{(۲) }\Big\|\tfrac{\partial \mathcal{L}}{\partial x_l}\Big\|=\Theta(1),\qquad
\text{(۳) }\|\Delta x_L\|_{\text{per step}}=\Theta(1) &&
\end{flalign*}
نرمال‌سازی این‌ها را به‌صورتِ پویا و در زمانِ اجرا برقرار می‌کند، اما این‌ها صرفاً شروطِ مقیاس‌اند. اگر همین شروط به‌صورتِ ایستا تضمین شوند---صفرکردنِ شاخه‌ها واریانسِ پیش‌رو را حفظ می‌کند (۱)، و کوچک‌سازیِ وابسته به عمقِ وزن‌ها نُرمِ گرادیان و اندازهٔ به‌روزرسانی را کران‌دار می‌کند (۲) و (۳)---آنگاه چشم‌اندازِ بهینه‌سازی به همان اندازه خوش‌وضع است و آموزش بدونِ هیچ نرمال‌سازیِ زمانِ‌اجرا پایدار می‌ماند. پس نرمال‌سازی شرطِ کافی برای آموزش‌پذیری است، نه لازم؛ آنچه لازم است حفظِ همین نامتغیّرهای مقیاس و خوش‌وضعی است که از راهِ مقداردهیِ اولیه و مقیاس‌بندیِ مناسبِ شاخهٔ باقی‌مانده نیز به‌دست می‌آید. تطبیقِ تجربیِ \lr{Fixup}/\lr{T-Fixup}/\lr{DeepNorm} با خطوطِ پایهٔ نرمال‌شده همین را تأیید می‌کند.